
        \documentclass[fleqn]{article}      
        \usepackage{amsmath}
        \usepackage{fontspec}
        \usepackage{kotex} % 한국어 지원  

        \begin{document}      
         쉬운 문제

\noindent 1. 주어진 행렬 \( A = \begin{pmatrix} 1 & 2 \\ 3 & 4 \end{pmatrix} \)의 행렬식(det)을 구하시오. \\\\
  
\noindent 2. 행렬 \( B = \begin{pmatrix} 5 & 6 \\ 7 & 8 \end{pmatrix} \)와 \( C = \begin{pmatrix} 1 & 0 \\ 0 & 1 \end{pmatrix} \)의 곱 \( BC \)를 구하시오. \\\\
  
\noindent 3. \( D = \begin{pmatrix} 2 & -3 \\ 4 & 1 \end{pmatrix} \)의 전치행렬 \( D^T \)를 구하시오. \\\\
  
\noindent 4. 다음 행렬의 크기를 구하시오: \( E = \begin{pmatrix} 1 & 2 & 3 \\ 4 & 5 & 6 \end{pmatrix} \). \\\\
  
\noindent 5. 행렬 \( F = \begin{pmatrix} 0 & 2 \\ -1 & 3 \end{pmatrix} \)의 역행렬 \( F^{-1} \)을 구하시오. \\\\

 중간 문제

\noindent 6. 다음 행렬의 합을 구하시오: \( G = \begin{pmatrix} 1 & 2 \\ 3 & 4 \end{pmatrix} \)와 \( H = \begin{pmatrix} 5 & 6 \\ 7 & 8 \end{pmatrix} \). \\\\
  
\noindent 7. 행렬 \( I = \begin{pmatrix} 2 & 0 \\ 1 & 3 \end{pmatrix} \)와 \( J = \begin{pmatrix} 1 & 4 \\ 2 & 5 \end{pmatrix} \)의 곱 \( IJ \)를 구하시오. \\\\
  
\noindent 8. \( K = \begin{pmatrix} 3 & -1 \\ 2 & 4 \end{pmatrix} \)의 대각합을 구하시오. \\\\
  
\noindent 9. \( L = \begin{pmatrix} 1 & 2 & 3 \\ 0 & 1 & 4 \\ 5 & 6 & 0 \end{pmatrix} \)의 자원행렬을 찾으시오. \\\\
  
\noindent 10. 행렬 \( M = \begin{pmatrix} 1 & 0 \\ 0 & 1 \end{pmatrix} \)에 대해, \( M^3 \)을 구하시오. \\\\

 어려운 문제

\noindent 11. 행렬 \( N = \begin{pmatrix} 1 & 2 & 3 \\ 0 & 1 & 4 \\ 0 & 0 & 1 \end{pmatrix} \)의 고유값을 구하시오. \\\\
  
\noindent 12. 행렬 \( O = \begin{pmatrix} 2 & 4 \\ 1 & 3 \end{pmatrix} \)의 고유벡터를 구하시오. \\\\
  
\noindent 13. 다음 연립방정식을 행렬로 표현하시오: \( 2x + 3y = 5 \), \( 4x - y = 1 \). \\\\
  
\noindent 14. 행렬 \( P = \begin{pmatrix} 1 & 2 \\ 3 & 4 \end{pmatrix} \)와 \( Q = \begin{pmatrix} 5 & 6 \\ 7 & 8 \end{pmatrix} \)의 고유값을 구하시오. \\\\
  
\noindent 15. 예를 들어 \( R = \begin{pmatrix} 1 & 2 \\ 2 & 4 \end{pmatrix} \)의 계수를 구하시고, 행렬의 선형 독립성을 판단하시오. \\\\

 정답

1. \( -2 \) \\\\\\\\
2. \( \begin{pmatrix} 5 & 6 \\ 7 & 8 \end{pmatrix} \) \\\\\\\\
3. \( \begin{pmatrix} 2 & 4 \\ -3 & 1 \end{pmatrix} \) \\\\\\\\
4. \( 2 \times 3 \) \\\\\\\\
5. \( \begin{pmatrix} \frac{3}{7} & -\frac{2}{7} \\ \frac{1}{7} & 0 \end{pmatrix} \) \\\\\\\\
6. \( \begin{pmatrix} 6 & 8 \\ 10 & 12 \end{pmatrix} \) \\\\\\\\
7. \( \begin{pmatrix} 5 & 28 \\ 7 & 19 \end{pmatrix} \) \\\\\\\\
8. \( 7 \) \\\\\\\\
9. \( 3 \) \\\\\\\\
10. \( M = \begin{pmatrix} 1 & 0 \\ 0 & 1 \end{pmatrix} \) \\\\\\\\
11. 고유값: \( 1, 1, 1 \) \\\\\\\\
12. 고유벡터: \( \begin{pmatrix} 2 \\ 1 \end{pmatrix} \) \\\\\\\\
13. \( \begin{pmatrix} 2 & 3 \\ 4 & -1 \end{pmatrix} \begin{pmatrix} x \\ y \end{pmatrix} = \begin{pmatrix} 5 \\ 1 \end{pmatrix} \) \\\\\\\\
14. 고유값: \( -1, 5 \) \\\\\\\\
15. 계수: \( 2 \); 선형 종속 \\\\\\\\      
        \end{document} 
    